% !TEX root = ./main.tex
In this paper we advocate a programming pattern for adaptations that
possibly reconfigure the architecture of service-oriented
applications.
%
Our pattern is supported by \SEArch, a platform featuring
semantic-based dynamic discovery and binding of services.
%
We identified some reconfigurations policies on the adaptable TeaStore
benchmark to demonstrate our approach; in particular, our policies fit
with the scenarios in~\cite[Sections 3.2.3 and
3.2.4]{bliudze:arXiv-2412.16060}.
%
This benchmark includes other scenarios that we did not consider here;
however, we argue that our approach is general enough to express the
other scenarios specified in~\cite{bliudze:arXiv-2412.16060} since the
architectural reconfigurations they encompass depend only on exogenous
control data.

As recalled in \cref{sec:search}, \SEArch contracts are formalised
in terms of extended CFSMs.
%
This allows us to specify contracts concerning application-level
(non-)functional requirements.
%
For instance, the quality of the image hinted in our running example
could be modelled as a quantitative attribute, say $Q$.
%
As mentioned in \cref{sec:progpat} we assume that control data are
exogenous.
%
This implies that they cannot include application-level information;
for instance, $Q$ cannot be used in the control loop to trigger an
adaptation when the image provided is sub-optimal (in fact, our
solution in \cref{sec:teapat} handles this in the application contract
of \cref{fig:external-service-contract} and not in adaptation
contracts).
%
Overcoming this limitation requires the introduction of mechanisms to
specify the interactions among the monitor and the execution
scenarios.
%
This can be readily supported in \SEArch since the monitor can be
conceived as a service local to a node whose behaviour is specified by
a suitable extended CFSM.
%
Even further, our platform could also support  monitors  as outsourced services
that could be dynamically discovered and used by other services once
suitable behavioural contracts are established for their usage.
%
These generalisations are in scope of future work.

Our programming pattern hinges on the control loop presented in
\cref{sec:progpat} that, liaising with a monitor that observes control
data, decides if a run-time architectural reconfiguration is necessary
and triggers the application-dependent activities provided by the
developer.
%
One could simplify this architecture by embedding the  control loop
in the monitor.
% 
We opted for a solution with two separated components because it
yields more flexiblility and can be extended more easily.
%
For instance, the generalisation to endogenous control data can be
attained more straightforwardly in our solution that in one where
control loop and monitor are entangled.

Our control loop can --at an abstract
level-- be envisaged as \quo{combinator} that composes the
application-dependent functions (the adaptation scenarios) necessary
for reconfigurations.
%
In this paper, we adopted a rather simplistic mechanism for the
composition of these functions.
%
Depending on the nature of the reconfigurations required by an
application, it could be necessary to appeal to more sophisticated 
composition mechanisms.
%
For instance, it could be necessary that the control loop passes
some data when invoking the adaptation scenario.
%
This can be supported by our platform at the cost of
more complex interactions between the control loop and execution
scenarios.

Finally, we recall that, for simplicity, the running adaptation
scenario is abruptly terminated when a reconfiguration is triggered
(line \ref{lst:finishing} in \cref{lst:control-loop}).
%
This may not be desirable; we plan to develop mechanisms for
guaranteeing graceful termination in future work.

\medskip

\noindent\textbf{Acknowledgements.}
We thank reviewers for their thorough reports and their suggestions that helped us to improve the paper.

% The key idea driving the design and implementation of \SEArch was to
% develop an execution architecture providing dynamic and transparent
% reconfiguration of service-based software artefacts.
% %
% Under \SEArch, services are brokered by a dedicated component capable
% of resolving compliance between requirements and provisions, by
% statically analysing formal interoperability contracts.
% %
% Noteworthy, contracts are formal and, besides specifying the expected
% communication protocol among participants, they can also declare
% functional and non-functional constraints.
% %
% For simplicity, we considered here only contracts specifying the
% communication protocols, but \SEArch can easily support adaptation
% scenarios using (non-)functional contracts as well.

% The possibility of describing each alternative execution as a
% different communication channel provides the support for considering
% alternative executions, hence adaptation scenarios, depending on the
% control data monitored from the environment.
% %
% Therefore, supporting adaptability as the possibility of choosing the
% best fitted implementation according to data monitored from the
% execution environment; presented in this work as a programming
% pattern.
% %
% Although not shown here, our programming patter can be used to
% implement in \SEArch also the other adaptation scenarios described
% in~\cite{bcglv12} once one has identified the relevant control data.


%%% Local Variables:
%%% mode: LaTeX
%%% TeX-master: "main"
%%% End:
